\section{Boundary Conditions}
\label{sec:bc}

Three boundary conditions are implemented, selected per family by the case file:
\texttt{farfield}, \texttt{slip\_wall}, and \texttt{no\_slip\_adiabatic\_wall}.

\subsection{Ghost state versus face state}
\label{sec:ghostvsface}

A design point worth stating explicitly, because it is a common source of silent
error, is that each boundary condition supplies \emph{two} different exterior
states and the solver uses them in different places:

\begin{itemize}
  \item The \textbf{ghost state} $\mathbf{U}^{\mathrm{ghost}}$ is a mirrored or
  characteristic state passed to the Riemann solver as the right state. Its
  purpose is that the resulting numerical flux \emph{enforces} the boundary
  condition.
  \item The \textbf{face state} $\mathbf{U}^{\mathrm{face}}$ is the physical
  state that actually exists on the boundary. It is used for the least-squares
  gradient stencil, for the limiter envelope, for the viscous wall flux, and for
  the surface output written to \texttt{surface.csv}.
\end{itemize}

Conflating the two is what produces the classic artefacts: doubled wall-normal
gradients, wall rows in the surface output that are really adjacent cell-centre
values, and a no-slip wall that reports a nonzero velocity. The distinction also
means \texttt{surface.csv} carries genuine boundary values, which is declared in
the metadata as \texttt{wall\_boundary\_output\_semantics = boundary\_value}.
Cell-centre quantities are written separately to
\texttt{surface\_cell\_center.csv} so the two can never be confused.

\subsection{Inviscid slip wall}

The ghost state reflects the normal momentum component while preserving density
and total energy,
\begin{equation}
  \mathbf{u}^{\mathrm{ghost}}=\mathbf{u}-2(\mathbf{u}\!\cdot\!\mathbf{n})\mathbf{n},
  \qquad
  \rho^{\mathrm{ghost}}=\rho,\qquad
  (\rho E)^{\mathrm{ghost}}=\rho E,
  \label{eq:slipghost}
\end{equation}
which is consistent because reflection leaves $|\mathbf{u}|$ unchanged. Feeding
\cref{eq:slipghost} to the Roe flux gives
\begin{equation}
  \mathbf{H}^{\mathrm{wall}}=\left[0,\ p\,n_x + \varepsilon n_x,\
  p\,n_y + \varepsilon n_y,\ 0\right]^T,
  \qquad \varepsilon \to 0 \ \text{ as } u_n\to0,
  \label{eq:slipflux}
\end{equation}
that is, \textbf{exactly} zero mass flux and \textbf{exactly} zero energy flux, with a
normal force that reduces to the pure pressure force $p\mathbf{n}$ only in the limit of
vanishing normal velocity.

An earlier draft of this report stated \cref{eq:slipflux} as an exact identity with
$\varepsilon\equiv0$, and that was wrong; the correction is recorded here rather than
quietly made because the claim had also been listed as a verified check. The mirrored
state has a normal-velocity jump of $\Delta u_n=-2u_n$, so the acoustic wave amplitudes
in the Roe dissipation do not vanish and their contributions do not cancel in the
normal-momentum component. Probing the solver's own \texttt{roeFlux} with the mirrored
state (via \texttt{report/check\_slip\_flux.sh}, which links the shipped library rather
than reimplementing it) gives, for $p=0.7143$:

\begin{center}
\begin{tabular}{lrrrr}
  \toprule
  Interior state, normal & Mass & Normal mom. & Energy & $u_n$ \\
  \midrule
  $\mathbf{V}=(0.3,0.1)$, $\mathbf{n}=+y$ & $0$ & $0.8244$ & $0$ & $0.100$ \\
  $\mathbf{V}=(0.9,-0.2)$, $\mathbf{n}=(0.6,0.8)$ & \num{-2.8e-17} & $1.2441$ & $0$ & $0.380$ \\
  $\mathbf{V}=(1,0)$, $\mathbf{n}=+y$ & \num{-5.4e-17} & $0.7143$ & \num{-1.7e-16} & $0$ \\
  $\mathbf{V}=0$ & $0$ & $0.7143$ & $0$ & $0$ \\
  \bottomrule
\end{tabular}
\end{center}

\noindent The two rows with $u_n=0$ --- a purely tangential flow and a stagnation point
--- return the pressure exactly. The rows with $u_n\ne0$ carry an excess that grows with
$u_n$. The mass and energy components are exactly zero in every case, which is the part
of the identity that matters for conservation: no mass or energy crosses a slip wall
regardless of the iterate.

The practical significance is small but should be stated rather than assumed. On a
converged solution the wall condition is satisfied, so $u_n$ at a slip wall is at
machine level --- the sanity checks record a maximum $|u_n|$ of
$\num{1.0e-13}$, $\num{1.4e-13}$ and $\num{1.9e-13}$ on the three inviscid cases --- and
the excess is correspondingly negligible in the reported forces. The discrepancy is a
property of \emph{intermediate} iterates, where it acts as a restoring term driving
$u_n$ towards zero, which is the intended behaviour of a ghost-state wall treatment.
What is \emph{not} justified is calling the identity exact, and it is no longer claimed
as such here or in \cref{sec:verification}.

The corresponding face state removes the normal velocity entirely,
$\mathbf{u}^{\mathrm{face}}=\mathbf{u}-(\mathbf{u}\!\cdot\!\mathbf{n})\mathbf{n}$,
keeps $\rho$ and $p$, and corrects the total energy for the reduced kinetic
energy. A slip wall therefore reports near-zero \emph{normal} velocity with a
retained tangential velocity, exactly as the output contract requires.

\subsection{Viscous no-slip adiabatic wall}

The ghost state reverses the entire velocity vector,
$\mathbf{u}^{\mathrm{ghost}}=-\mathbf{u}$, with $\rho$ and $\rho E$ unchanged, so
that the arithmetic face average of the interior and ghost velocities is exactly
zero. The face state imposes the wall condition directly: zero velocity,
interior pressure ($\partial p/\partial n=0$), and interior temperature
(adiabatic, $\partial T/\partial n=0$), with the density following from the
equation of state,
\begin{equation}
  \mathbf{u}^{\mathrm{face}}=\mathbf{0},\quad
  p^{\mathrm{face}}=p_i,\quad
  T^{\mathrm{face}}=T_i,\quad
  \rho^{\mathrm{face}}=\frac{p^{\mathrm{face}}}{R\,T^{\mathrm{face}}},\quad
  (\rho e)^{\mathrm{face}}=\frac{p^{\mathrm{face}}}{\gamma-1}.
  \label{eq:noslipface}
\end{equation}

Two points of honesty about \cref{eq:noslipface}. First, with $p^{\mathrm{face}}=p_i$
and $T^{\mathrm{face}}=T_i$ the equation of state returns the interior density
bit-for-bit, so writing $\rho=p/(RT)$ is a statement of \emph{why} the wall density
equals the interior density rather than a computation that produces a different value;
the code says as much in a comment at the same line. It is the correct form because it
is what makes the adiabatic assumption explicit and would give the right answer if
either the pressure or the temperature condition were changed, but no claim of
sophistication over a copy is intended. Second, the total energy is written here as
$\rho e$ rather than $\rho E$ because the velocity is zero: the general definition of
\cref{sec:equations} includes the kinetic term, which vanishes at a no-slip wall, and
using the same symbol for the reduced quantity would be inconsistent with the rest of
the report.

Because the surface output uses this state, the reported wall velocity and Mach
number are identically zero to machine precision. The adiabatic condition is
additionally enforced on the flux by projecting the normal component out of
$\nabla T$, as described after \cref{eq:wallgrad}.

\subsection{Farfield}

The farfield is characteristic, based on locally one-dimensional Riemann
invariants along the face normal. With interior values $(\rho_i,p_i,u_{n,i},a_i)$
and freestream values $(\rho_\infty,p_\infty,u_{n,\infty},a_\infty)$, the normal
Mach number $M_n=u_{n,i}/a_i$ selects the branch:

\begin{itemize}
  \item $M_n\ge 1$ (supersonic outflow): the boundary state is the interior state;
  nothing is imposed.
  \item $M_n\le -1$ (supersonic inflow): the boundary state is the freestream.
  \item $|M_n|<1$ (subsonic): the outgoing and incoming invariants
  \begin{equation}
    R^+=u_{n,i}+\frac{2a_i}{\gamma-1},\qquad
    R^-=u_{n,\infty}-\frac{2a_\infty}{\gamma-1}
  \end{equation}
  are combined into
  $u_{n,b}=\tfrac12(R^++R^-)$ and $a_b=\tfrac14(\gamma-1)(R^+-R^-)$. Entropy
  $s=p/\rho^\gamma$ and the tangential velocity are taken from the upwind side
  (interior if $u_{n,b}>0$, freestream otherwise), giving
  $\rho_b=\left(a_b^2/(\gamma s)\right)^{1/(\gamma-1)}$ and
  $p_b=s\rho_b^\gamma$.
\end{itemize}

This treatment is what allows the supersonic cases to be posed on the same mesh
as the subsonic ones without special handling: at $\Minf=2$ the inflow arc is
fully determined by the freestream and the outflow arc is fully determined by the
interior, so the bow shock can leave the domain cleanly. Degenerate invariant
combinations that would give $a_b\le0$, or a non-positive reconstructed density or
pressure, fall back to the freestream state rather than producing an unphysical
boundary value. The same characteristic state is used as both ghost and face
state at the farfield, since no wall condition applies there.

\subsection{Force integration and the pressure/friction split}
\label{sec:forces}

Forces are integrated over wall faces only. Recalling from \cref{sec:mesh} that
the stored boundary normal $\mathbf{n}$ points out of the fluid and into the
body, the outward normal of the \emph{body} surface is $\mathbf{m}=-\mathbf{n}$,
and the Cauchy traction exerted by the fluid on the body is
\begin{equation}
  \mathbf{t}=\boldsymbol{\sigma}\!\cdot\!\mathbf{m}
  =\left(-p\mathbf{I}+\boldsymbol{\tau}\right)\!\cdot\!(-\mathbf{n})
  =p\,\mathbf{n}-\boldsymbol{\tau}\!\cdot\!\mathbf{n}.
  \label{eq:traction}
\end{equation}
The pressure force per face is therefore $+(p_{\mathrm{wall}}-p_\infty)\mathbf{n}\,\Delta s$
--- referenced to $p_\infty$ so that a uniform ambient pressure contributes
nothing on a closed body --- and the viscous force is $-\boldsymbol{\tau}\!\cdot\!\mathbf{n}\,\Delta s$.
Getting this sign pair right is essential: with the opposite viscous sign a flat
plate with $\partial u/\partial y>0$ would produce thrust instead of drag.

The wall pressure is reconstructed from the adjacent cell with the same gradient
used by the flux, so the integrand is second-order accurate and consistent with
the residual. For the friction the \emph{tangential} part of the viscous
traction is used,
\begin{equation}
  \boldsymbol{\tau}_t=\boldsymbol{\tau}\!\cdot\!\mathbf{n}
  -\left[(\boldsymbol{\tau}\!\cdot\!\mathbf{n})\!\cdot\!\mathbf{n}\right]\mathbf{n},
  \label{eq:shear}
\end{equation}
which is what the contract requires: the normal viscous traction is explicitly
excluded rather than being mislabelled as skin friction. The wall-normal velocity
derivatives entering $\boldsymbol{\tau}$ are corrected with the imposed zero wall
velocity exactly as in \cref{eq:wallgrad}.

Each wall face is integrated by the rank that owns its adjacent cell, so a face on
a partition boundary is counted exactly once globally, and the six accumulators
(pressure force, viscous force, moment, wall length) are combined in a single
\texttt{MPI\_Allreduce}. The moment about the case reference centre is
$M_z=\sum (\mathbf{r}\times\mathbf{f})_z$ with
$\mathbf{r}=\mathbf{x}_f-\mathbf{x}_{\mathrm{ref}}$, positive counter-clockwise.
